\documentclass[11pt]{article}

\usepackage[margin=1in]{geometry}
\usepackage{times}
\usepackage{booktabs}
\usepackage{longtable}
\usepackage{array}
\usepackage{xcolor}
\usepackage{listings}
\usepackage{hyperref}
\usepackage{enumitem}
\usepackage{titlesec}
\usepackage{parskip}

\lstset{
    basicstyle=\ttfamily\small,
    breaklines=true,
    frame=single,
    backgroundcolor=\color{gray!8},
    columns=fullflexible,
    xleftmargin=0.5em,
    xrightmargin=0.5em
}

\hypersetup{
    colorlinks=true,
    linkcolor=blue,
    urlcolor=blue
}

\title{PendingIntent Vulnerability Class (v1/v2/v3):\\
Static Triage, Cross-App Exploitation, and Verified Mitigation}
\author{Research / Test Evaluation Report}
\date{August 24, 2026}

\begin{document}

\maketitle

\begin{abstract}
This report documents a full evaluation cycle for a well-known Android vulnerability
class involving \texttt{PendingIntent} objects: unrestricted mutability (fill-in Intent
injection), privilege escalation via delegated \texttt{send()} capability, and ownership
hijacking via \texttt{cancel()}. The evaluation proceeds through four stages of
increasing rigor: (1) static triage over a 245-APK F-Droid/IzzyOnDroid corpus, (2) a
same-process Frida-based simulation of attack and defense logic, (3) a genuine
two-application, two-UID exploitation of the vulnerability built and executed on a real
Android emulator, and (4) a genuine two-application mitigation, likewise built and
executed, using OS-enforced controls rather than in-app logic. The final result is
empirical, not projected: two of three attack variants were confirmed exploitable in an
unprotected baseline and confirmed blocked in a hardened variant, on-device.
\end{abstract}

\tableofcontents

\section{Background and Objective}

\texttt{PendingIntent} is an Android API object that lets one application delegate a
piece of its own authority (the ability to fire an \texttt{Intent} with that
application's identity) to another component or process, most commonly the system
server (for notifications, alarms, and app widgets). Because a \texttt{PendingIntent}
is a Binder-backed token, whoever holds a reference to it can, subject to platform
protections, invoke operations on it as if they were the creator. Three attack
variants were investigated:

\begin{itemize}[leftmargin=1.5em]
    \item \textbf{v1 -- Fill-in Intent / mutable PendingIntent.} If a
    \texttt{PendingIntent} is created without \texttt{FLAG\_IMMUTABLE}, a holder can
    supply a ``fill-in'' \texttt{Intent} via \texttt{send(Context, int, Intent)} that
    overrides the action, data, component, or categories of the underlying Intent,
    redirecting what the creator's privileged operation actually does.
    \item \textbf{v2 -- Privilege / permission boundary.} A \texttt{PendingIntent}
    handed to an external component (an exported activity/service/receiver, or leaked
    via an unprotected broadcast) may let an unrelated app trigger the creator's
    privileged action via \texttt{send()}.
    \item \textbf{v3 -- PendingIntent cancel / ownership.} A holder that is not the
    creator may call \texttt{cancel()} on a captured \texttt{PendingIntent}, silently
    revoking the creator's own capability.
    \end{itemize}

The objective of this evaluation was to determine, with evidence rather than assumption,
whether these three variants are exploitable across a process/UID boundary, and whether
a proposed set of mitigations genuinely prevents them at the operating-system level.

\section{Evaluation Strategy}

The evaluation was carried out in four escalating stages, each addressing a limitation
identified in the previous one.

\subsection*{Stage 1 -- Static corpus triage}
A Python analyzer (\texttt{analyze\_pendingintent\_corpus.py}, built on
\texttt{androguard 4.x}) scanned 245 real-world APKs drawn from the \texttt{fdroid\_main}
and \texttt{izzyondroid} mirrors. For every method invoking a
\texttt{PendingIntent.*} factory or operation, the analyzer inspected the raw decoded
bytecode for:
\begin{itemize}[leftmargin=1.5em]
    \item Folded mutability-flag bit patterns (\texttt{FLAG\_MUTABLE},
    \texttt{FLAG\_IMMUTABLE}, \texttt{FLAG\_ALLOW\_UNSAFE\_IMPLICIT\_INTENT}), detected
    by bitwise testing of every integer literal operand rather than exact string
    matching, since the compiler folds flag combinations into a single constant.
    \item String/API indicators for fill-in usage, external Intent sourcing
    (\texttt{getParcelableExtra}, etc.), ownership checks
    (\texttt{getCreatorPackage}, \texttt{getCallingUid}, etc.), and permission checks.
    \item Manifest-level exposure: exported components, permission-guarded components,
    declared permissions.
\end{itemize}
Each APK received a v1, v2, and v3 heuristic score with an explicit disclaimer in the
tool's own documentation: \emph{``This is STATIC TRIAGE. A v1/v2/v3 score does NOT prove
exploitability.''} Thirteen APKs (5 v1, 5 v2, 3 v3) were selected for downstream dynamic
validation.

\subsection*{Stage 2 -- Same-process Frida simulation}
A Frida harness (\texttt{run\_validation.py} plus \texttt{pendingintent\_attack\_sim.js}
/ \texttt{pendingintent\_defense\_sim.js}) attached to each of the 13 selected APKs,
hooked the \texttt{PendingIntent} factory methods, and on every capture attempted three
attacker-style operations (fill-in injection, \texttt{send()}, \texttt{cancel()}) either
unconditionally (attack mode) or behind a reference ownership-authorization gate
(defense mode). This stage's own documentation flags its central limitation: because
Frida attaches to and executes inside the target app's own process, there is no genuine
second UID involved -- the ``attacker'' identity is a hardcoded, synthetic string
compared against the real \texttt{creatorPackage} the OS reports. Running this against
the 13 real candidates also revealed a practical limitation: most of them never created
an observable \texttt{PendingIntent} during the passive observation window (see Table
\ref{tab:v3-defense}), so no exploit or defense outcome was actually recorded for those
apps.

\subsection*{Stage 3 -- Genuine cross-app, cross-UID exploitation}
To remove the same-process limitation of Stage 2, two purpose-built Android applications
were designed, implemented, compiled, signed, and installed as if they were real,
independent apps: \texttt{VictimApp}, which creates and leaks a
\texttt{PendingIntent}, and \texttt{MalwareApp}, a genuinely separate application (own
package, own UID, own process) that captures the leaked \texttt{PendingIntent} and
performs the three attacker operations against the live Android API. This is a real
two-app reproduction of the attack, not a simulation of one.

\subsection*{Stage 4 -- Genuine cross-app mitigation and verification}
A hardened variant of each app (\texttt{VictimAppHardened}, \texttt{MalwareAppV2}) was
built to test whether OS-enforced controls -- rather than any application-level logic --
actually prevent the attacks. \texttt{MalwareAppV2} was deliberately signed with a
different certificate than \texttt{VictimAppHardened}, which is the precondition for a
signature-level permission check to have anything to enforce. Both the baseline and the
hardened scenario were then executed end-to-end on a real Android emulator, and the
logcat output was captured as direct evidence.

\section{Why an In-Receiver Ownership Check Cannot Work}

Before designing Stage 4's mitigation, an initially tempting approach was considered and
rejected: adding an ownership check inside the receiver that handles the
\texttt{PendingIntent}'s underlying action (e.g., checking
\texttt{Binder.getCallingUid()} inside \texttt{onReceive()}). This does not work,
because when Android delivers a broadcast triggered via \texttt{PendingIntent.send()},
the receiving component is invoked with no reliable signal identifying which process
called \texttt{.send()}. The entire purpose of \texttt{PendingIntent} is delegation: the
action executes with the \emph{creator's} authority regardless of who invoked it, and
\texttt{Binder.getCallingUid()} inside the receiver typically reflects the broadcast
dispatcher (e.g. \texttt{system\_server}), not the actual invoking app. Any check
written at that point would be guessing, not enforcement. Real prevention must instead
be applied \emph{upstream}, at the point where the \texttt{PendingIntent} is created and
shared, using controls the OS itself enforces.

\section{Application Details}

Four Android applications were designed and implemented as minimal, buildable, real
APKs (no Gradle; built directly with SDK command-line tools). All target
\texttt{minSdkVersion=23}, \texttt{targetSdkVersion=30}, matching the only system image
available in the local environment (Pixel 6 AVD, \texttt{android-30/google\_apis/x86\_64}).

\begin{table}[h]
\centering
\begin{tabular}{@{}p{2.6cm}p{3.1cm}p{3.6cm}p{4.3cm}@{}}
\toprule
\textbf{App} & \textbf{Package} & \textbf{Role} & \textbf{Signing Certificate} \\
\midrule
VictimApp & \texttt{com.research.pivictim} & Creates a \emph{mutable} PendingIntent
(\texttt{FLAG\_MUTABLE}), leaks it via an unprotected exported broadcast
(\texttt{PI\_LEAK\_ACTION}). & \texttt{debug.keystore}, CN=Android Debug,
SHA-256 \texttt{ec26c9ed\ldots} \\
MalwareApp & \texttt{com.research.pimalware} & Captures the leaked PendingIntent
(receiver registered both statically and dynamically) and fires attacks A/B/C against
the real Android API. & Same as VictimApp (same cert; UID differs) \\
VictimAppHardened & \texttt{com.research.pivictimhardened} & Creates an
\emph{immutable} PendingIntent (\texttt{FLAG\_IMMUTABLE}); leak broadcast requires a
newly declared \texttt{protectionLevel="signature"} permission. & Same
\texttt{debug.keystore} as VictimApp \\
MalwareAppV2 & \texttt{com.research.pimalwarev2} & Declares
\texttt{uses-permission} for the hardened victim's signature permission; identical
attack logic, retargeted. & \texttt{malware\_v2.keystore}, CN=Malware V2 Attacker,
SHA-256 \texttt{8504e94f\ldots} (deliberately different) \\
\bottomrule
\end{tabular}
\caption{The four test applications built for Stages 3 and 4.}
\end{table}

\subsection*{VictimApp / VictimAppHardened -- Sensitive Operation}
Both variants target the same conceptual ``privileged operation'':
\texttt{SensitiveActionReceiver}, a broadcast receiver representing whatever sensitive
action a real app might gate behind a \texttt{PendingIntent} (dismissing a
notification, completing a payment, replaying an authenticated action). It has no
ownership check of its own -- by design, since Section 3 establishes that such a check
cannot work at that point -- and every invocation is logged to Logcat and to an
internal marker file, whether triggered by the app's own logic or by an external
\texttt{PendingIntent.send()}/\texttt{cancel()} call.

\subsection*{MalwareApp / MalwareAppV2 -- Attack Runner}
Both variants implement identical attack logic in an \texttt{AttackRunner} class:

\begin{lstlisting}[language=Java, caption={Attack B, privilege escalation via bare send()}]
private static void attackB_privilegeEscalation(PendingIntent captured) {
    try {
        captured.send();
        Log.i(TAG, "{\"attack\":\"B_privilege_escalation\",\"status\":\"EXECUTED\"}");
    } catch (SecurityException e) {
        Log.i(TAG, "{...\"status\":\"SECURITY_EXCEPTION\"...}");
    }
}
\end{lstlisting}

A build-time correction was required during implementation: Android 8+ (API 26+)
background execution limits silently drop most broadcasts to a manifest-declared
receiver once the receiving app is not in the foreground -- which occurs here because
launching the victim's activity backgrounds the malware app moments later. Both
malware apps were updated to \emph{additionally} register their capture receiver
dynamically via \texttt{Context.registerReceiver()}, a technique real malware commonly
uses for the same reason. This change does not weaken the Stage 4 test: permission
checks are enforced by \texttt{BroadcastQueue} at delivery time regardless of how the
receiver was registered, and this was directly confirmed in the captured logs (Section
6.2).

\section{Build and Verification Methodology}

All four APKs were built entirely with Android SDK command-line tools (no Gradle):

\begin{enumerate}[leftmargin=1.5em]
    \item \textbf{Compile:} \texttt{javac}, JDK 26.0.2, against
    \texttt{android-30/android.jar}.
    \item \textbf{Dex:} \texttt{d8} (build-tools 36.0.0).
    \item \textbf{Resource compile/link:} \texttt{aapt2 compile} then \texttt{aapt2 link}.
    \item \textbf{Dex injection:} \texttt{jar --update} to insert \texttt{classes.dex}
    into the linked base APK.
    \item \textbf{Alignment:} \texttt{zipalign -p 4}.
    \item \textbf{Signing:} \texttt{apksigner sign}, using two distinct keystores as
    described in Table 1.
\end{enumerate}

Prior to any runtime test, static correctness was verified directly:

\begin{itemize}[leftmargin=1.5em]
    \item \texttt{apksigner verify --print-certs} confirmed the two signing
    certificates differ (Table 1), the precondition for a signature-level permission
    check to be meaningful at all.
    \item \texttt{aapt dump xmltree AndroidManifest.xml} on the compiled
    \texttt{pivictimhardened.apk} confirmed the declared permission's
    \texttt{protectionLevel} compiles to the binary value \texttt{0x2} (signature),
    not merely appearing correct in source XML.
\end{itemize}

Runtime verification used a Pixel 6 AVD (API 30, x86\_64, Google APIs, 1536MB RAM),
booted via \texttt{emulator.exe} with WHPX hardware acceleration, confirmed reachable
via \texttt{adb devices} and \texttt{adb shell getprop sys.boot\_completed}.

\section{Results}

\subsection{Static Triage (Stage 1)}

Of the 13 selected candidates, the three intended as ``pure v3'' (cancel/ownership
risk) were in fact all classified \texttt{mixed} by the scoring logic -- meaning each
also scored $\geq 10$ on at least one other axis. No app in the 245-APK corpus was
found to exhibit \emph{isolated} v3 risk; real-world PendingIntent cancel/ownership
exposure co-occurred with fill-in/mutability or permission-boundary exposure in every
case observed.

\begin{table}[h]
\centering
\small
\begin{tabular}{@{}lrrrl@{}}
\toprule
\textbf{Package} & \textbf{v1} & \textbf{v2} & \textbf{v3} & \textbf{Classification} \\
\midrule
com.adguard.android.contentblocker & 19 & 18 & 23 & mixed \\
org.secuso.privacyfriendly2048 & 15 & 18 & 23 & mixed \\
com.newoether.agora & 15 & 18 & 23 & mixed \\
\bottomrule
\end{tabular}
\caption{Static v1/v2/v3 scores for the 3 selected v3 candidates.}
\end{table}

\subsection{Same-Process Frida Simulation (Stage 2)}

\label{tab:v3-defense}
\begin{table}[h]
\centering
\small
\begin{tabular}{@{}lrlll@{}}
\toprule
\textbf{Package} & \textbf{PI captured} & \textbf{Attack A} & \textbf{Attack B} & \textbf{Attack C} \\
\midrule
com.adguard.android.contentblocker & 1 & NULL\_PENDING\_INTENT & NULL\_PENDING\_INTENT & NULL\_PENDING\_INTENT \\
org.secuso.privacyfriendly2048 & 0 & NO\_DATA & NO\_DATA & NO\_DATA \\
com.newoether.agora & 1 & NULL\_PENDING\_INTENT & NULL\_PENDING\_INTENT & NULL\_PENDING\_INTENT \\
\bottomrule
\end{tabular}
\caption{Defense-mode Frida results for the 3 real v3 candidates: no real,
non-null PendingIntent was ever captured during the observation window.}
\end{table}

No exploit or defense outcome was empirically recorded for the real F-Droid candidates
at this stage; the passive observation window never coincided with a real
\texttt{PendingIntent} creation event for these apps.

\subsection{Cross-App Baseline Exploitation (Stage 3) -- Real, Verified}

Executed on the Pixel 6 / API 30 emulator. Full captured Logcat output (tags
\texttt{PIVictim}, \texttt{PIMalware}):

\begin{lstlisting}
I/PIVictim: Created mutable PendingIntent: ... creatorPackage=com.research.pivictim creatorUid=10167
I/PIVictim: Leaked PendingIntent via broadcast action=com.research.pivictim.PI_LEAK_ACTION
I/PIMalware: PI_CAPTURED pendingIntent=... creatorPackage=com.research.pivictim
             creatorUid=10167 thisPackage=com.research.pimalware thisUid=10168
I/PIMalware: {"attack":"A_fillin_injection","status":"EXECUTED"}
I/PIMalware: {"attack":"B_privilege_escalation","status":"EXECUTED"}
I/PIMalware: {"attack":"C_cancel_hijack","status":"SECURITY_EXCEPTION",
             "detail":"java.lang.SecurityException: Permission Denial:
             cancelIntentSender() from pid=4873, uid=10168 is not allowed
             to cancel package com.research.pivictim"}
I/PIVictim: SensitiveActionReceiver fired: EXECUTED ... attackExtra=A_fillin_injection_from_pimalware
I/PIVictim: SensitiveActionReceiver fired: EXECUTED ... attackExtra=null
\end{lstlisting}

\textbf{Unexpected finding.} Attack C (cancel-ownership hijack) was \emph{not}
exploitable even in this fully unprotected baseline. Android's own
\texttt{ActivityManagerService.cancelIntentSender()} independently checks the calling
UID against the PendingIntent's owning package and throws \texttt{SecurityException} on
mismatch, regardless of mutability flags or broadcast permissions. This corrected an
initial assumption in the evaluation plan (that Attack C would succeed unprotected).
Attacks A and B, by contrast, both executed successfully and were confirmed by
\texttt{VictimApp}'s own receiver logging the attacker-controlled data.

\subsection{Cross-App Hardened Mitigation (Stage 4) -- Real, Verified}

Executed in the same emulator session. Full captured Logcat output (tags
\texttt{PIVictimHardened}, \texttt{BroadcastQueue}; note the complete absence of any
\texttt{PIMalwareV2} line):

\begin{lstlisting}
I/PIVictimHardened: Created IMMUTABLE PendingIntent: ...
                     creatorPackage=com.research.pivictimhardened creatorUid=10169
I/PIVictimHardened: Sent signature-permission-gated broadcast
                     action=com.research.pivictimhardened.PI_LEAK_ACTION
                     requiredPermission=com.research.pivictimhardened.permission.RECEIVE_PI_LEAK
W/BroadcastQueue: Permission Denial: receiving Intent { act=...PI_LEAK_ACTION ... }
                  to ProcessRecord{... com.research.pimalwarev2/u0a170} (pid=4956, uid=10170)
                  requires com.research.pivictimhardened.permission.RECEIVE_PI_LEAK
                  due to sender com.research.pivictimhardened (uid 10169)
W/BroadcastQueue: Permission Denial: receiving Intent { act=...PI_LEAK_ACTION ... }
                  to com.research.pimalwarev2/.PiCaptureReceiver
                  requires com.research.pivictimhardened.permission.RECEIVE_PI_LEAK
                  due to sender com.research.pivictimhardened (uid 10169)
\end{lstlisting}

The Android OS itself (\texttt{BroadcastQueue}) denied delivery to
\texttt{com.research.pimalwarev2} (uid 10170, signed with a different certificate) at
the permission check, before \texttt{onReceive()} was invoked for either the
manifest-declared or the dynamically-registered receiver. No \texttt{PendingIntent} was
ever captured by \texttt{MalwareAppV2}; consequently no attack of any kind had a target
to operate on. \texttt{VictimAppHardened}'s own \texttt{SensitiveActionReceiver} log
shows only the creation and leak-attempt lines, confirming it never received an
attacker-triggered invocation.

\subsection{Consolidated Result}

\begin{table}[h]
\centering
\begin{tabular}{@{}lll@{}}
\toprule
\textbf{Attack} & \textbf{Baseline (unprotected)} & \textbf{Hardened} \\
\midrule
A -- fill-in injection & EXECUTED & Never reached (OS denied broadcast delivery) \\
B -- send() escalation & EXECUTED & Never reached (OS denied broadcast delivery) \\
C -- cancel() hijack & SECURITY\_EXCEPTION (platform's own UID check, unrelated to tested mitigations) & N/A -- no capture occurred \\
\bottomrule
\end{tabular}
\caption{Final, empirically observed outcome for all three attack variants under both scenarios.}
\end{table}

\section{Conclusions}

\begin{enumerate}[leftmargin=1.5em]
    \item \textbf{The vulnerability is real, not merely a static-analysis score.} Attacks
    A and B were demonstrated as genuine, working exploits between two independently
    signed, independently processed Android applications -- not a same-process
    simulation and not a heuristic score.
    \item \textbf{In-app ownership checks placed at the point of use (the receiver)
    cannot work}, because \texttt{PendingIntent} delegation removes any reliable signal
    of the true invoking process at that point.
    \item \textbf{Effective prevention must be applied upstream}, at PendingIntent
    creation (\texttt{FLAG\_IMMUTABLE}) and at the point where the PendingIntent is
    shared (a signature-level permission on the delivery channel). Both mechanisms are
    enforced by the Android platform itself, not by application logic, and both were
    empirically confirmed -- via direct OS-level denial in Logcat, not by inference --
    to stop the two attacks that were otherwise exploitable.
    \item \textbf{Attack C (cancel hijack) is independently mitigated by the platform}
    at API level 30 via UID validation inside \texttt{cancelIntentSender()}, regardless
    of the mitigations under test. This was an empirical correction to an initial
    assumption, not a predicted result.
    \item \textbf{Open scope for further work:} this evaluation confirms the
    vulnerability and its mitigation on synthetic fixture apps at a single API level
    (30). It does not establish that any of the 13 real F-Droid/IzzyOnDroid candidates
    identified by static triage are themselves exploitable, since none of them yielded
    an observable \texttt{PendingIntent} capture during passive dynamic observation
    (Table \ref{tab:v3-defense}). It also does not establish whether platform behavior
    (particularly the \texttt{cancelIntentSender()} UID check) is consistent across the
    full range of API levels present in the corpus (up to API 37).
\end{enumerate}

\end{document}
